\addcontentsline{toc}{chapter}{Conclusion}
\chapter*{Conclusion}

The current \textit{open AI} debate produces false binaries, paradoxes, and policy-failure. This research instead preoccupies itself with the \textit{democratisation} of AI, in an attempt to re-centre the intended societal-level outcomes of AI governance, rather than clashing over prior beliefs imported from related fields of study. 

`Democratising AI' involves establishing legitimate mechanisms for supporting equitable, inclusive and sustainable distributions of benefit and access around AI development, use of AI, value-generated by AI, and governance over AI systems themselves. This research evaluates DIs as a well-theorised institutionalised form of commons data governance, according to the extent to which they might be extended to facilitate AI's democratisation.

To make the case for DIs, I first unpack datasets as strategic and underappreciated policy objects within AI governance, and propose that establishing dedicated DIs to design, develop and steward datasets for ML uses (both training and benchmark) would insert a decentralised, democratic layer of community-accountability early in the ML value chain. This upstream positioning is important. Downstream system-output level accountability and decision points tend towards restrictive or distributive policy choices, whereas DIs' upstream location embeds grassroots empowerment of data-generating communities, alongside integrative creation of public ‘value’ (a deliberately expansive term, encompassing a range of possible dimensions for a range of possible stakeholders). 

To operationalise public `value', I chose to focus specifically on data quality, safe and responsible data usage, data-sustainability as well as equity and inclusion (other formulations would have been valid). Doing so, I showcase how DIs might address current concerns and open questions in data-centric AI governance, and to suggest how DIs might obtain buy-in from stakeholders across the policy-making process. This latter is important for designing robust and feasible policy instruments. Finally, I propose nine high-level principles for policy-makers to consider when designing DIs for ML datasets.

This research is intended to open up a new, complementary strain of AI policy thinking around dataset development, and to add a new lever for policy-makers to consider. This is not a definitive account. There are further dimensions of value to be evaluated, feasibility and implications details to be ironed out, and constraints to be attended to. It is a first pass at synthesising the knowledge of multiple relevant domains into a coherent agenda from which to extend and expand both DI and AI governance theorising.